\begin{frame}{Contexto da prática}
  \begin{block}{Problema investigado}
    Como a troca dos valores de dois resistores altera o ganho, o amortecimento e
    a resposta transitória de um filtro ativo passa-baixa de segunda ordem?
  \end{block}

  \begin{columns}[T,onlytextwidth]
    \begin{column}{0.48\textwidth}
      \textbf{Regime permanente senoidal}
      \begin{itemize}
        \item magnitude e fase de $H(j\omega)$;
        \item ganho em baixas frequências;
        \item frequência de corte de $-3\,\mathrm{dB}$;
        \item possível pico de ressonância.
      \end{itemize}
    \end{column}
    \begin{column}{0.48\textwidth}
      \textbf{Regime transitório}
      \begin{itemize}
        \item resposta a uma entrada em degrau;
        \item valor final da tensão de saída;
        \item tempos $t_{10}$ e $t_{90}$;
        \item resposta monótona ou oscilatória.
      \end{itemize}
    \end{column}
  \end{columns}
\end{frame}
